\chapter{Research Methodology}
The adopted research methodology comprises three main phases, namely preparation of datasets for conducting experiments, applying the proposed and existing methods to do experiments, and evaluation of different methods.

\setcounter{equation}{0}

\section{Data Collection}
Table~\ref{tab:dsets} provides the properties of the text corpora considered to evaluate feature selection methods. Documents with multiple class labels are not considered for clustering. The variation in the class distribution of a dataset is expressed using the coefficient of variation \cite{ram2020}. The skewness in the class distribution of a dataset is expressed using the adjusted Fisher-Pearson standardized moment coefficient of skewness \cite{weltman2018},\cite{kovchegov2020}. All classes of Ohsumed dataset have exactly same number of documents (400), so it's coefficient of skewness is not defined (ND). The Sparsity ratio is defined as $1-(nze/n \cdot m)$, where $nze$ is the number of non zero elements in the VSM representation of a dataset.


\begin{center}
\begin{table}
\begin{threeparttable}
\caption{Properties of Datasets}
\label{tab:dsets}
\begin{tabular}{lrrrrrr}
\toprule
\multirow{2}{*}{Dataset} & \multicolumn{3}{c}{Number of} & \multicolumn{2}{c}{Coefficient of} & \multirow{2}{*}{Sparsity ratio}\\
			\cline{2-6}
             & documents & features & classes & variation & skewness &   \\
\midrule
Reuters\tnote{1}      &  8213	 & 18933 & 	41   & 	3.23  & 4.81  & 0.997  \\     
Ohsumed\tnote{2}     & 	9200 	 & 13510 &  23   & 	0  	  &  ND   & 0.996  \\     
Newsgroups\tnote{1}   & 	18846	 & 26214 & 	20	 &  0.10  &-2.39  & 0.997  \\     
Ohscal\tnote{2}       &  11162	 & 11465 &  10	 &  0.27  & 0.24  & 0.995  \\     
Sports\tnote{2}       &  8580	 & 14870 & 	 7	 &  1.02  & 1.01  & 0.991  \\     
Hitech\tnote{2}       &  2301	 & 12940 &   6	 &  0.50  &-0.62  & 0.989  \\     
Reviews\tnote{2}    	 &  4069	 & 22925 & 	 5	 &  0.64  &-0.43  & 0.992  \\     
Classic\tnote{2}     &  7095	 & 7747 & 	 4	 &  0.55  & 1.77  & 0.995  \\     
Webkb\tnote{2}        &  8282	 & 22890 & 	 7	 &  1.06  & 1.74  & 0.996  \\     
\bottomrule
\end{tabular}
\begin{tablenotes}
\item[1] http://www.becbapatla.ac.in/TextData.html
\item[2] http://www.textcollections.com
\end{tablenotes}
\end{threeparttable}
\end{table}
\end{center}




\section{Clustering}
\subsection{Motivation for the selection of K-means}
\begin{inparaenum}[\itshape i\upshape)]
\item The time complexities of clustering algorithms are given in \cite{xu2015}. K-means algorithm is scalable because its time complexity is $\mathcal{O}(nkum)$. Where $n$ is the number of documents in the text collection, $k$ is the number of classes in the dataset, $u$ is the terminating condition which is the maximum number of iterations allowed and $m$ is the number of dimensions used to represent the text dataset. 
\item The algorithm is easy to implement. 
\item The algorithm has known limitations. 
\item The parameters of the algorithm can be fine tuned easily to improve the performance \cite{franti2018}. The algorithm is vulnerable to outliers because, in addition to its sensitivity to the initial centroids, it can fall in the local regions in
the search space. To overcome this the algorithm is run multiple times on the dataset with different initial centroids and best performance is reported.

\end{inparaenum}
 
\subsection{Parameters used for K-means}

Initialization method and repetition of K-means algorithm has significant impact on the performance of the clustering algorithm.
The parameters used for K-means algorithm are shown in Table~\ref{tab:pkmeans}. 

\renewcommand\arraystretch{2} 
\begin{table} [!h]
\centering
\caption{Parameters chosen for K-means clustering algorithm}
\begin{tabular}{m{7cm}m{7cm}}
\toprule
Parameter & Choice \\
\midrule
Distance measure, in m-dimensional space, that the algorithm should minimize & cosine \\

Method used to choose initial cluster centroid positions & Perform preliminary clustering phase on random 10\% subsample of dataset. Clustering in this phase is by randomly choosing $k$ objects from the subsample. \\

Number of times the algorithm is repeated, each with a new set of initial centroids & 10 \\

Action to take if any cluster loses all of its member objects & Create a new cluster consisting of the one object furthest from its centroid \\

Maximum number of iterations allowed for optimizing the objective function & 100 \\

\bottomrule
\end{tabular}
\label{tab:pkmeans}
\end{table}

\section{Algorithm in \LaTeX}
\begin{algorithm}
\begin{algorithmic}[1]
\REQUIRE Input array $A$
\ENSURE Sorted array $A$
\WHILE{($x<10$)}
   \STATE $x \gets x-1$
\ENDWHILE

\STATE Initialize variables
\FOR{each element $x$ in $A$}
    \IF{$x$ is out of order}
        \STATE Swap $x$ with previous element
    \ENDIF
\ENDFOR
\RETURN $A$
\end{algorithmic}
\caption{Algorithm example}
\label{algo:bs}
\end{algorithm}

\section{Code Snippets}

Here is a Python program:

\begin{minted}{python}
def is_prime(n):
    if n <= 1:
        return False
    for i in range(2, int(n**0.5) + 1):
        if n % i == 0:
            return False
    return True

print(is_prime(17))
\end{minted}